\documentclass{article}

\begin{document}

\section*{Differentiability}

\begin{itemize}
      \item{} \emph{Statement and proof of theorem 41 (Chain Rule for functions from $R^n$ to $R$)}

            \quad{}
            Theorem 41 allows us to take directional derivatives when giving a multivariable function, a base point, and a direction vector from that base point.
            It requires that all partial derivatives to be continuous because it describes the rate of change of the function when all basis vectors are changing with respect to a single variable t.
            Also since it involves a dot product if one of the partial derivatives is nonexistent than there would be no result.
            In a geometric sense, the gradient represents the maximum amount of change on the function, so the direction is the amount of projection that lies on that.
            In the proof it splits two points on the curve so you can analyze each rate of change by a single variable.
            Then it uses MVT to show that there is a rate of change for each basis direction and how it relates to the actual direction vector.
            The final chain rule is used to get the rate f changes with respect to t.

            \vspace{0.5cm}

      \item{} \emph{Statement and proof of theorem 42 (Differentiability and partial derivatives)}

            \quad{}
            The vector proposed as a in the differentiability definiiton is the gradient vector.
            The theorem is basically saying that the rate that the difference between f and the tangent
            plane created with regard to the gradient vector is going to zero faster than the rate the inputs
            are going to zero. Thereby proving that the partial derivatives have to be continuous in order for
            the whole function to be differentiable. It proves this by first subtracting the tangent plane approximation
            from the function. Then it bounds this "error" by utilizing that the partial derivatives are continuous.
            Thus, resulting in the error being bounded by epsilon since there exists a delta y, proving differentiability.
\end{itemize}

\end{document}